\uuid{xQXk}
\exo7id{7342}
\titre{exo7 7342}
\auteur{mourougane}
\organisation{exo7}
\datecreate{2021-08-10}
\isIndication{false}
\isCorrection{false}
\chapitre{Groupe, anneau, corps}
\sousChapitre{Anneau}
\cpge{mpsi-mp2i.polynomes-fractions-rationnelles.c, pcsi.polynomes.c, ptsi.polynomes.c, mpsi-mp2i.nombres-reels-suites.a, pcsi.raisonnement-vocabulaire-ensembliste.c, ptsi.raisonnement-vocabulaire-ensembliste.c, bcpst1.nombres-reels, bcpst1.polynomes-reels.b}
\module{Algèbre}
\niveau{L2}
\difficulte{}

\contenu{
\texte{
Le but de l'exercice est de déterminer une valeur approchée de la racine cubique réelle $c$ de $17$.
On considère le polynôme $P(x)=x^3-17$.
}
\begin{enumerate}
    \item \question{Montrer que $2<c<3$.}
    \item \question{Déterminer par la méthode de Hörner le polynôme $R(z):=10^3P(2+\frac{z}{10})$ en la variable $z$.}
    \item \question{Déterminer une valeur approchée de $c$ à $10^{-1}$ près.}
    \item \question{Déterminer une valeur approchée de $c$ à $10^{-2}$ près.
%(réponse : $2,57<c<2,58$.)
}
\end{enumerate}
}
